\scanpage{9}
\begin{flushright}
\handwritten{16 Juillet .}\\
\end{flushright}
Cher Grothendieck,\\

Voici la nouvelle rédaction de \handwritten{4}.\handwritten{10}.\handwritten{6}.

En relisant le n\textsuperscript{o} \handwritten{9} de S.G.A \handwritten{4} IV j'ai trouvé les broutilles suivantes :

p. \handwritten{149} l. \handwritten{16} lire \(E'\handwritten{\vee} E'' = E\)

p. \handwritten{158} l. \handwritten{5} lire \(s\unsure{\varepsilon|\in}\operatorname{Hom}(e_E,G)\)

p. \handwritten{158} compléter le diagramme \handwritten{9}.\handwritten{3}.\handwritten{6}.\handwritten{2}

p. \handwritten{159} l. \handwritten{15} lire ``le foncteur \(\handwritten{\sigma}\) composé de

Dans les n\textsuperscript{o} \handwritten{10} et suivant j'ai trouvé

p. \handwritten{2} l. -\handwritten{7} lire ``doù l'assertion ''

p. \handwritten{14} l. \handwritten{5} lire \(\handwritten{\mathscr{H}}\mathrm{om}_A\)

p.\handwritten{14} l. \handwritten{17} itou .

\margin{\handwritten{\figure{\begin{tikzpicture}[x=0.007cm,y=-0.007cm]
\node[anchor=base west,rotate=36,inner sep=0pt] at (110,1108) {\uline{demander}};
\node[anchor=base west,rotate=36,inner sep=0pt] at (137,1153) {\uline{Messing}};
\draw (325,1037) -- (323,1148) -- (326,1165);
\end{tikzpicture}}}}
Enfin , désagréable surprise , je me suis aperçu qu'il manquait le n\unsure{\deleted{?}\handwritten{\textsuperscript{o}}|\handwritten{\textsuperscript{o}}} \handwritten{14} dans les notes que tu m'as envoyés (Modules sur un topos défini par recollement ). Je me souviens avoir rédigé ce numéro et je pense bien te l'avoir donné . En tout cas je ne l'ai pas dans mes papiers . J'espère que ce numéro ne s'est pas perdu .

\begin{flushright}
Bien à toi\\
\handwritten{\unsure{J. L.|J. V.}}\\
\end{flushright}
% archive sheet number: 8
